\hypertarget{appendix-c-the-pep-hall-of-fame}{%
\chapter{Appendix C: The PEP Hall of
Fame}\label{appendix-c-the-pep-hall-of-fame}}

The history of Python is the history of its \textbf{Python Enhancement
Proposals (PEPs)}.

\begin{longtable}[]{@{}
  >{\raggedright\arraybackslash}p{(\columnwidth - 4\tabcolsep) * \real{0.33}}
  >{\raggedright\arraybackslash}p{(\columnwidth - 4\tabcolsep) * \real{0.33}}
  >{\raggedright\arraybackslash}p{(\columnwidth - 4\tabcolsep) * \real{0.33}}@{}}
\toprule
PEP \# & Title & Impact \\
\midrule
\endhead
\textbf{PEP 8} & Style Guide for Python Code & The standard for
readable, idiomatic Python. \\
\textbf{PEP 20} & The Zen of Python & The guiding philosophy of the
language. \\
\textbf{PEP 257} & Docstring Conventions & Formalized internal
documentation. \\
\textbf{PEP 343} & The ``with'' Statement & Introduced context managers
and resource safety. \\
\textbf{PEP 380} & Syntax for Delegating to a Subgenerator & Introduced
\passthrough{\lstinline!yield from!}. \\
\textbf{PEP 443} & Single-dispatch generic functions & Functional-style
polymorphism. \\
\textbf{PEP 484} & Type Hints & The foundation of modern Python static
typing. \\
\textbf{PEP 498} & Literal String Interpolation & Introduced
F-Strings. \\
\textbf{PEP 525} & Asynchronous Generators & Bridged the gap between
\passthrough{\lstinline!asyncio!} and
\passthrough{\lstinline!yield!}. \\
\textbf{PEP 572} & Assignment Expressions & The Walrus Operator
(\passthrough{\lstinline!:=!}). \\
\textbf{PEP 594} & Removing dead batteries & Cleaned up the Standard
Library for Python 3.13. \\
\textbf{PEP 634} & Structural Pattern Matching & Introduced
\passthrough{\lstinline!match!} and \passthrough{\lstinline!case!}. \\
\textbf{PEP 703} & Making the GIL Optional & The roadmap for
Free-Threaded Python. \\
\bottomrule
\end{longtable}

\begin{center}\rule{0.5\linewidth}{0.5pt}\end{center}

\textbf{END OF APPENDICES}

\begin{center}\rule{0.5\linewidth}{0.5pt}\end{center}

To achieve the ultimate level of ``Godhood,'' one must look beyond the
virtual machine and understand how Python interacts with physical
hardware.

\hypertarget{chapter-70-cpu-cache-locality-and-data-alignment}{%
\chapter{Chapter 70: CPU Cache Locality and Data
Alignment}\label{chapter-70-cpu-cache-locality-and-data-alignment}}

Modern CPUs are significantly faster than system memory. Performance is
often bottlenecked by the ``Memory Wall.''

\hypertarget{the-cache-hierarchy-l1-l2-l3}{%
\subsection{70.1 The Cache Hierarchy (L1, L2,
L3)}\label{the-cache-hierarchy-l1-l2-l3}}

When the CPU needs data, it checks the caches first. A cache hit takes
\textasciitilde1-10 cycles, while a main memory access (cache miss)
takes \textasciitilde200-300 cycles.

\hypertarget{why-python-is-cache-unfriendly}{%
\subsubsection{1. Why Python is
Cache-Unfriendly}\label{why-python-is-cache-unfriendly}}

Standard Python objects are scattered across the heap. A
\passthrough{\lstinline!list!} of \passthrough{\lstinline!float!}
objects is actually an array of pointers to
\passthrough{\lstinline!PyObject!} structs. * \textbf{Pointer Chasing}:
To read the value of \passthrough{\lstinline!mylist[0]!}, the CPU must
load the pointer, then jump to another memory location to load the
actual float value. This jump often causes a cache miss.

\hypertarget{the-solution-array.array-and-numpy}{%
\subsubsection{\texorpdfstring{2. The Solution: \texttt{array.array} and
NumPy}{2. The Solution: array.array and NumPy}}\label{the-solution-array.array-and-numpy}}

As seen in Chapter 28, contiguous memory is the secret. By storing raw
C-types in a block, the CPU can pre-fetch the next values into the
cache, leading to 10x-100x speedups for numerical processing.

\hypertarget{memory-alignment-and-padding}{%
\subsection{70.2 Memory Alignment and
Padding}\label{memory-alignment-and-padding}}

C-structs (like those in Chapter 24) are padded by the compiler to
ensure that fields start at memory addresses divisible by their size
(e.g., an 8-byte double should start at an 8-byte boundary). *
\textbf{Performance}: Misaligned access can require two memory fetches
instead of one, or even trigger hardware exceptions on some
architectures.

\begin{center}\rule{0.5\linewidth}{0.5pt}\end{center}
